\documentclass[10pt,a4paper]{amsart}
\usepackage[margin=19mm]{geometry}
\usepackage{amsmath,amssymb,mathtools}
\usepackage[scr=rsfs,cal=euler]{mathalfa}
\usepackage{xfrac}
\usepackage[unicode,hidelinks,bookmarks=false]{hyperref}
\newcommand{\ie}{i.e.\ }
\newcommand{\cf}{cf.\ }
\newcommand{\resp}{respectively\ }
\newcommand{\Subsection}{\subsection}
\providecommand{\nobreakdash}{\nobreak\hbox{-}}
\setlength{\emergencystretch}{2em}
\multlinegap0pt
\usepackage{amssymb}
\usepackage{dsfont}
\usepackage{enumerate}
\usepackage{xcolor}
\usepackage{bm}
\usepackage{tikz}
\newtheorem{theorem}{Theorem}
\newtheorem{lemma}{Lemma}
\usepackage{cleveref}
\crefname{theorem}{Theorem}{Theorems}
\crefname{lemma}{Lemma}{Lemmas}
\newcommand{\Cay}{\operatorname{Cay}}
\newcommand{\OO}{\mathcal O}
\title{Even-Intersecting Families of Permutations}
\author{Anirban Banerjee, Abisek Dewan, Rajiv Mishra}
\date{Source: arXiv:2609.21645v1 (CC BY 4.0)}
\begin{document}
\maketitle

\noindent\emph{Source and licence.} Even-Intersecting Families of Permutations. Version arXiv:2609.21645v1 (CC BY 4.0), distributed by arXiv under the Creative Commons Attribution 4.0 International licence, http://creativecommons.org/licenses/by/4.0/, as recorded in the arXiv licence metadata for this exact version and shown on the abstract page https://arxiv.org/abs/2609.21645v1. Full licence notice: \texttt{licenses/arxiv-2609.21645-CC-BY-4.0.txt}.\par\smallskip

\noindent\textit{Editorial scope.} Counted unit: the article's theorem thm:main\_result (source0.tex lines 222--232). The article's complete original proof of it is delivered with it (source0.tex lines 1204--1309, one \texttt{proof} environment of 106 lines, which the article places away from the statement). No mathematics is written, repaired or re-derived: the statement and the proof are the article's own lines, in the article's own notation, and no retained line is substituted or re-flowed.  Retained uncounted as the source-local context the proof needs: lines 570--573; lines 996--1017.  Nothing was omitted from any retained span, so the package carries no omission record.  No retained line contains a citation, so the wrapper carries no bibliography and no bibliography entry.  No retained line contains a figure environment, an image-inclusion command or drawing code, so no image is delivered and the package declares no asset dependency.  Editorial formatting only, in addition: an AMS \texttt{amsart} preamble that restates the article's own declarations and macros used by the retained lines, namely the environments \texttt{theorem}, \texttt{lemma} on separate counters, together with the standard packages those lines need.  The retained environments print in this document as Theorem 1, Lemma 1 in order of appearance, whereas the article prints the same items with the numbering of its own full document.  The complete native source of the article as distributed in the arXiv e-print for this exact version is bundled unchanged as \texttt{sources/210976/source0.tex}.
\begin{lemma}\label{lem:Weighted_Adjacency_eigenvlues}
 Let $X_1,X_2,...,X_t$ be conjugacy classes of $S_n$. For real numbers $c_1,c_2,...,c_t$ define $A=\sum_{i=1}^{t}c_i A_i$, where $A_i$ is the adjacency matrix of the normal Cayley graph $\Cay(S_n,X_i)$. If $g_i \in X_i$, then the eigenvalue of $A$ corresponding to irreducible character $\chi^{\rho}$ is given by
    \[\theta_{\rho}=\frac{1}{\dim(\rho)} \sum_{i=1}^{t}c_i|X_i| \chi^{\rho}(g_i)\]
\end{lemma}

\begin{theorem}\label{thm:z_upper_bound}
Let 
$
\mathcal{O}_n = \{ \rho \vdash n : m_1(\rho) \text{ is odd} \} 
$ be a subset of partitions with odd number of fixed points.
Suppose $\{x_\lambda\}$ is a set of non-negative weights indexed by partitions $\lambda \vdash n$ with $\lambda \neq (n)$, such that
$$
\sum_{\substack{\lambda \vdash n \\ \lambda \neq (n)}} x_\lambda = 1.
$$
Furthermore, assume that for every $\rho \in \mathcal{O}_n$, a real number $z$ satisfies the bound
$$
z \le -\sum_{\substack{\lambda \vdash n \\ \lambda \neq (n)}} x_\lambda \frac{\chi^\lambda(\rho)}{f^\lambda}.
$$
Then, there exist absolute constants $C>0$ and $n_0$ such that for all even $n \ge n_0$,
$$
z \le \frac{\exp{\left(\frac{n}{2}+Cn^{3/4}\log n\right)}}{(n-1)!!}.
$$
Consequently, for sufficiently large $n$
$$
z \le \frac{e^{\frac{n}{2}+o(n)}}{(n-1)!!}.
$$
\end{theorem}

\begin{theorem}\label{thm:main_result}
Let $M(n)$ be the maximum size of an even-intersecting family of $S_n$.
    Then, there exist absolute constants
$C>0$ and $n_0 \in \mathbb{N}$ such that, for every even $n\geq n_0$,
\begin{equation*}
n!!\leq M(n)
\leq e^{{n}/{2}+C n^{3/4}\log n}\,n!!,
\label{eq:main-even-intro}
\end{equation*}
where the lower bound is attained by $\mathcal{F}=S_2\wr S_{n/2}$.
\end{theorem}

\begin{proof}[Proof of $\Cref{thm:main_result}$]
Let
$$
\mathcal{O}_n
:=
\{\rho\vdash n:m_1(\rho)\text{ is odd}\},
$$
and, for $\rho\in\mathcal{O}_n$, let $X_\rho$ denote the conjugacy class of
$S_n$ of cycle type $\rho$. 
Let $A_\rho$ be the adjacency matrix of $\Cay(S_n,X_\rho)$, for each $\rho\in \OO_n$ and consider the weighted adjacency matrix
$$
A
:=
\sum_{\rho\in\mathcal{O}_n}c_\rho A_\rho.
$$
By \Cref{lem:Weighted_Adjacency_eigenvlues}, for every
$\lambda\vdash n$, the corresponding eigenvalue is
$$
\theta_\lambda
=
\sum_{\rho\in\mathcal{O}_n}
c_\rho |X_\rho|\frac{\chi^\lambda(\rho)}{f^\lambda}.
$$

We want to find the weights $c_\rho$, $\rho\in \OO_n$ such that $\theta_{(n)}=1$ is the largest eigenvalue and the smallest eigenvalue is as large as possible. This leads to the following linear program:

\noindent\textbf{Primal LP:}

\noindent Minimize $y$ subject to:
\begin{eqnarray*}
    \sum_{\rho \in \OO} \left( \frac{\chi^{\lambda}(\rho) }{f^{\lambda}} \right) \omega_{\rho} + y &\ge& 0, \quad \text{for all } \lambda \neq (n),\\
    \sum_{\rho \in \OO} \omega_{\rho} &=& 1,\\
    \omega_\rho&\geq& 0, \qquad\rho\in\OO.
\end{eqnarray*}

In the above linear program, $\omega_\rho=c_\rho|X_\rho|$. We restrict to nonnegative weights $\omega_\rho$, and hence to $c_\rho\geq 0$. Although signed weights are allowed in the weighted Hoffman framework, this restriction merely narrows the class of admissible matrices and still yields a valid upper bound for $M(n)$.

Feasibility of the primal LP is easy: Take $\omega_{(n-1,1)}=1$ and all other $\omega_\rho$ zero, and $y=2$. Determining the optimal weights
$\omega_\rho$ for the primal LP directly appears difficult, therefore we consider the dual of this LP as follows.


\noindent\textbf{Dual LP:}

\noindent Maximize $z$ subject to:
\begin{eqnarray*}
    \sum_{\substack{\lambda\vdash n\\\lambda\neq(n)}} \left( \frac{\chi^{\lambda}(\rho) }{f^{\lambda}} \right) x_{\lambda} + z &\leq& 0, \quad \text{for all } \rho \in \OO,\\
    \sum_{\substack{\lambda\vdash n\\\lambda\neq(n)}} x_{\lambda} &=& 1,\\
    x_{\lambda}&\geq& 0, \qquad \lambda\vdash n,\; \lambda\neq (n).
\end{eqnarray*}
First, we construct a feasible solution for the dual. 
Take  
\begin{align*}
    z&=\frac{1}{(n-1)!!-1},\\
    x_{2\nu}&=\frac{f^{2\nu}}{(n-1)!!-1}\qquad \text{for all }\nu\vdash n/2,\;\nu\neq (n/2),
\end{align*}
and put all other $x_\lambda=0$. It is easy to verify that this is a feasible solution for every $n\geq 6$. In fact, we have verified in the SageMath for $n=6$ to $16$ that this is indeed an optimal solution. Assuming this as an optimal solution, we get $\theta_{\min}=-1/((n-1)!!-1)$, and therefore, Hoffman's ratio bound will give
$$M(n)\leq \frac{n!}{1+(n-1)!!-1}=n!!,$$
which will answer the Cameron--Deza--Frankl problem, but proving the optimality is difficult. 






Let $z_n$ be the optimal solution for the dual. The purpose of the
character optimization problem studied in the previous section was precisely
to obtain a sufficiently strong upper bound for $z_n$. By
\Cref{thm:z_upper_bound}, there exist absolute constants
$C>0$ and $n_0$ such that, for every even $n\geq n_0$, we have
$$
z_n
\leq
\frac{
\exp\left(
\frac{n}{2}
+
C n^{3/4}\log n
\right)
}{
(n-1)!!
}.
$$
Let $y_n$ be the optimal solution of the primal, then by strong duality, we have
$$
y_n=z_n,
$$
and hence there exists a choice of weights $\omega_\rho$ for which the least
eigenvalue of the corresponding weighted adjacency matrix satisfies
$$\theta_{\min}\geq-\frac{\exp\left(\frac{n}{2}+C n^{3/4}\log n\right)}{(n-1)!!}.$$


Applying the weighted Hoffman bound gives
\begin{align*}
    M(n)&\leq\frac{n!}{1-\frac{1}{\theta_{\min}}}\\
    &\leq n!!\exp\left(\frac{n}{2}+C n^{3/4}\log n\right).
\end{align*}

Equivalently, for sufficiently large even $n$
$$
M(n)\leq e^{\, n/2+o(n)}\,n!!.
$$

For the lower bound, let $n=2m$. In its natural action on $[n]$, the wreath product $H=S_2\wr S_m\le S_n$ consists of all permutations preserving the partition into pairs
$$B_i=\{2i-1,2i\},\qquad 1\le i\le m.$$
Thus an element of $H$ may permute the pairs and independently swap the two points within each pair. Therefore it gives an even-intersecting family of size $|S_2\wr S_{n/2}|=n!!$.
\end{proof}

\end{document}
